% Source: Weinberg GR, physical PDF p. 218 (printed p. 195).
% Coverage: Figure 8.2, geometric elements of an eccentric ellipse.
% Status: vector reconstruction; source-reviewed and compile-clean.
% Uncertainties: none.

\begingroup
\renewcommand{\thefigure}{8.2}
\begin{figure}[t]
\centering
\begin{tikzpicture}[x=0.82cm,y=0.82cm,>=Latex]
  \draw[line width=0.75pt] (0,0) ellipse (4.4 and 2.25);
  \coordinate (F) at (3.35,0);
  \fill (F) circle (1.8pt);

  % Aphelion and perihelion distances measured from the occupied focus.
  \draw[<->] (-4.4,0) -- (F);
  \node[above] at (-0.55,0) {$r_+$};
  \draw[<->] (F) -- (4.4,0);
  \node[below] at (3.88,-0.02) {$r_-$};

  % Semilatus rectum at the focus.
  \draw[<->] (F) -- (3.35,1.45);
  \node[left] at (3.28,0.83) {$L$};

  % Semimajor axis and focal offset above the ellipse.
  \draw (0,2.45) -- (0,4.18);
  \draw (4.4,2.28) -- (4.4,4.18);
  \draw[<->] (0,3.92) -- (4.4,3.92);
  \node[above] at (2.2,3.92) {$a$};

  \draw (3.35,2.10) -- (3.35,3.35);
  \draw[<->] (0,3.10) -- (3.35,3.10);
  \node[above] at (1.68,3.10) {$ea$};

  % Semiminor axis.
  \draw (4.95,0) -- (5.95,0);
  \draw (0,2.25) -- (5.95,2.25);
  \draw[<->] (5.65,0) -- (5.65,2.25);
  \node[right] at (5.75,1.15) {$a\sqrt{1-e^2}$};
\end{tikzpicture}
\caption{Elements of an ellipse referred to in the calculation of the
precession of planetary orbits. (The ellipse here has the same eccentricity
as the orbit of Icarus.)}
\label{fig:8.2}
\end{figure}
\endgroup
